\documentclass[english,landscape,xcolor=dvipsnames,aspectratio=169]{beamer}
% aspectratio= Other possible values are: 1610, 149, 54, 43 and 32. Default is 43

%\usepackage[brazil]{babel}
%\usepackage[utf8]{inputenc}
\usepackage{ragged2e} %justifying
\usepackage{caption}
%\usepackage{subcaption}
\usepackage{graphicx}
\usepackage{subfig}
\usepackage{tikz}
\usepackage{pgfplots}
\usetikzlibrary{spy,calc}
\usepackage{hyperref}
\usepackage{listings}
\usepackage{tikzpagenodes}
\usepackage{UERJpresentation}


%New Commands
\newcommand{\Bi}{\mathrm{Bi}}
\newcommand{\nmax}{n_{\max}}
\newcommand{\dd}{\mathrm{d}}



\newcommand{\event}{COBEM}
\newcommand{\eventdate}{December 2017}
\newcommand{\mytitle}{Template Example: Solution of the Heat Conduction in Solid-State Electronics by Integral Transforms}



\author[\scriptsize 
Prof. Daniel J. N. M. Chalhub -- daniel.chalhub@uerj.br
]
{
\textbf{Prof. Daniel J. N. M. Chalhub} (Presenter)
\\ 	
Lívia M. Corrêa
}

\institute[UERJ]{

	\begin{minipage}{\textwidth}
		\begin{columns}
			\column{1\textwidth}
			\centering
			Dep. of Mech. Engineering -- PPG-EM,\\
			Universidade do Estado do Rio de Janeiro, \\
			Rio de Janeiro - RJ, Brazil
		\end{columns}
	\end{minipage}
	
	
	\begin{minipage}{\textwidth}
		\vspace{1em}
		\begin{columns}
		    \column{0.15\textwidth}
			\column{0.15\textwidth}
			\centering
			\includegraphics[height=3.5em]{Logo-UERJ.png}
			\column{0.2\textwidth}
			\centering
			\includegraphics[height=4em]{LogoPPGEM-color.png}
			\column{0.2\textwidth}
			\centering
			\includegraphics[height=3.0em]{Logo-GESAR.png}
			\column{0.1\textwidth}
		\end{columns}
	\end{minipage}
	
	\begin{minipage}{\textwidth}
	\vspace{1em}
		\begin{columns}
			\column{1\textwidth}
			\centering
			\normalsize \emph{\event{} -- \eventdate}
		\end{columns}
	\end{minipage}
			
}






\begin{document}

%%%%%%%%%%%%%%%%%%%%%%%%%%%%
%%%%%%%%%%%%%%%%%%%%%%%%%%%%
\begin{frame}
    \titlepage
\end{frame}
%^^^^^^^^^^^^^^^^^^^^^^^^^^
%^^^^^^^^^^^^^^^^^^^^^^^^^^



%~~~~~~~~~~~~~~~~~~~~~~~~~~
\section{Introduction}
%~~~~~~~~~~~~~~~~~~~~~~~~~~




	
%%%%%%%%%%%%%%%%%%%%%%%%%%%%
%%%%%%%%%%%%%%%%%%%%%%%%%%%%	
\begin{frame}{Introduction}\includelogo
One of the biggest challenges on the design of modern electronic devices is the thermal control.
		
\vspace{1em}
		
In this work, it is proposed a solution by Classical Integral Transform Technique(CITT) to solve a heat conduction problem on Solid-State Electronics (SSE) considering heat generation and external convection.
		
\end{frame}
%^^^^^^^^^^^^^^^^^^^^^^^^^^
%^^^^^^^^^^^^^^^^^^^^^^^^^^







%%%%%%%%%%%%%%%%%%%%%%%%%%%%
%%%%%%%%%%%%%%%%%%%%%%%%%%%%
\begin{frame}{Solid State Electronics (SSE)}\includelogo

	
Solid-state electronics (SSE) are those circuits or devices built entirely from solid materials and in which the electrons, or other charge carriers, are encapsulated within the solid material.
\vspace{0.5em}	    
	            
Exemples: Transistors, Microprocessors, Random Access Memory (RAM), Itegrated Circuits (IC). 
	        
\begin{figure}
	\centering
	%\caption{Contour plot of the heat generation considered cases.} \label{heatgenimage}
	\hspace*{-0.5cm}
	\subfloat{%: $\xi_0=1/2$ and $\eta_0=1/2$]{%\label{centralheat}
		\fbox{\includegraphics[width=0.2\hsize]{Figuras/ic.jpg}}}
	\hspace*{0.2cm}
	\subfloat{%  $\xi_0=1/4$ and $\eta_0=2/5$]{%\label{initialheat}
		\fbox{\includegraphics[width=0.25\hsize]{Figuras/ram.jpg}}}
	\hspace*{0.2cm}
	\subfloat{%  $\xi_0=1/4$ and $\eta_0=2/5$]{%\label{initialheat}
		\fbox{\includegraphics[width=0.25\hsize]{Figuras/processor.jpg}}}
\end{figure}
\end{frame}
%^^^^^^^^^^^^^^^^^^^^^^^^^^
%^^^^^^^^^^^^^^^^^^^^^^^^^^












%%%%%%%%%%%%%%%%%%%%%%%%%%%%
%%%%%%%%%%%%%%%%%%%%%%%%%%%%
\begin{frame}{Substrate}\includelogo


Substrate is a solid substance and serve as the foundation upon which electronic devices such as transistors, diodes, and especially Integrated Circuits (ICs) are encapsulated.

\vspace{0.5em}	

Example of substrate's materials:

\begin{itemize}
\item silicon 
\item silicon dioxide 
\item aluminum oxide 
\item sapphire 
\end{itemize} 


The advantage of this is the superior insulation between adjacent transistors.

\end{frame}
%^^^^^^^^^^^^^^^^^^^^^^^^^^
%^^^^^^^^^^^^^^^^^^^^^^^^^^









%%%%%%%%%%%%%%%%%%%%%%%%%%%%
%%%%%%%%%%%%%%%%%%%%%%%%%%%%
\begin{frame}{Exemple of SSE: Integrated Circuit}\includelogo
\begin{figure}
\centering
\subfloat{
	\fbox{\includegraphics[width=0.4\hsize]{Figuras/microchip.jpg}}}
\hspace*{0.2cm}
\subfloat{
	\fbox{\includegraphics[width=0.3\hsize]{Figuras/microchip-microscopio.jpg}}}

\end{figure}
\end{frame}
%^^^^^^^^^^^^^^^^^^^^^^^^^^
%^^^^^^^^^^^^^^^^^^^^^^^^^^







%~~~~~~~~~~~~~~~~~~~~~~~~~~
\section{Mathematical Formulation}
%~~~~~~~~~~~~~~~~~~~~~~~~~~



%%%%%%%%%%%%%%%%%%%%%%%%%%%%
%%%%%%%%%%%%%%%%%%%%%%%%%%%%
\begin{frame}{Problem Description}\includelogo
\centering\includegraphics[width=0.8\textwidth]{Figuras/volcontrol}	
\end{frame}
%^^^^^^^^^^^^^^^^^^^^^^^^^^
%^^^^^^^^^^^^^^^^^^^^^^^^^^





	
	
	
	
%%%%%%%%%%%%%%%%%%%%%%%%%%%%
%%%%%%%%%%%%%%%%%%%%%%%%%%%%	
\begin{frame}{Problem Formulation}\includelogo
The nondimensionalization of the problem leads to the following mathematical formulation:
%
\small
\begin{gather*}\label{equation}
\frac{\partial^2\Theta}{\partial \xi^2}+
\beta^2 \frac{\partial^2\Theta}{\partial \eta^2}-(\Bi \gamma) \Theta  =-G(\xi,\eta)  \quad \text{for} \quad 0\leq \xi \leq 1 \quad \text{and} \quad  0\leq \eta \leq 1
\\
\left.\frac{\partial\Theta}{\partial \xi}\right|_{\xi=0}=0;	\qquad		\left.\frac{\partial\Theta}{\partial \eta}\right|_{\eta=0}=0; \\ \left.\frac{\partial\Theta}{\partial \xi}\right|_{\xi=1}=0;	\qquad \left.\frac{\partial\Theta}{\partial \eta}\right|_{\eta=1}=0;
\end{gather*}
\end{frame}
%^^^^^^^^^^^^^^^^^^^^^^^^^^
%^^^^^^^^^^^^^^^^^^^^^^^^^^






	
%~~~~~~~~~~~~~~~~~~~~~~~~~~
\section{Methodology}
%~~~~~~~~~~~~~~~~~~~~~~~~~~






%%%%%%%%%%%%%%%%%%%%%%%%%%%%
%%%%%%%%%%%%%%%%%%%%%%%%%%%%
\begin{frame}{Classical Integral Transform Technique}\includelogo
In order to solve the proposed problem, the Classical Integral Transform Technique (CITT) is applied, considering the appropriate Helmholtz Eigenvalue Problem in cartesian coordinates.
\begin{gather*}
\Psi_n''(\eta)+\lambda_n^2\Psi_n(\eta)=0 \\
\Psi_n'(0)=0; \quad \Psi_n'(1)=0.
\end{gather*}
where $\Psi(\eta)$ are the eigenfunctions and $\lambda_n$ are the eigenvalues. For this particular problem, the case where $\lambda=0$ also exists.

For $\lambda=0$, the solution of the eigenvalue problem is given by:
\begin{gather*}
\Psi_0(\eta)=1; \qquad \lambda_0=0;
\end{gather*} 
%
And for $\lambda>0$:
\begin{gather*}
\Psi_n(\eta)=\cos(\lambda_n \eta); \qquad \lambda_n=n \pi, \quad \text{for} \quad n=1,2,3,\dots
\end{gather*} 
\end{frame}
%^^^^^^^^^^^^^^^^^^^^^^^^^^
%^^^^^^^^^^^^^^^^^^^^^^^^^^




	
	
	
	
	
	
	
%%%%%%%%%%%%%%%%%%%%%%%%%%%%
%%%%%%%%%%%%%%%%%%%%%%%%%%%%	
\begin{frame}{Classical Integral Transform Technique}\includelogo
To apply the CITT, the transformation pair is defined.
\begin{varblock}[10cm]{Transformation Pair}
\begin{align*}
%\label{transformada}\quad
\text{Transformation} \quad &\Rightarrow \quad 
\bar{\Theta}_n(\xi)=\int_{0}^{1}\Theta\Psi_n(\eta)\dd\eta\\
%\qquad\label{eq:inversion}
\text{Inversion} \quad & \Rightarrow \quad \Theta= \sum_{n=0}^{\infty} \frac{ \bar{\Theta}_n(\xi)\Psi_n(\eta)}{N_n} 
\end{align*}
\vspace{0em}
\end{varblock}
$N_n$ is the norm and is defined as:
%
\begin{gather*}
N_n=\int_{0}^{1}\Psi_n^2 \dd \eta  
\end{gather*}
\end{frame}
%^^^^^^^^^^^^^^^^^^^^^^^^^^
%^^^^^^^^^^^^^^^^^^^^^^^^^^




	
	
	
%%%%%%%%%%%%%%%%%%%%%%%%%%%%
%%%%%%%%%%%%%%%%%%%%%%%%%%%%
\begin{frame}{Classical Integral Transform Technique}\includelogo
The dimensionless diferential equation is written again, \\ 
multiplied by $\Psi_n$ and integrated in the domain for $\eta$. 
\begin{gather*}
	\int_{0}^{1}\frac{\partial ^2 \Theta}{\partial\xi ^2} \Psi_n d\eta+\beta^2\int_{0}^{1}\frac{\partial ^2 \Theta}{\partial\eta ^2}\Psi_n d\eta - \Bi\gamma\int_{0}^{1}\Theta\Psi_n d\eta=- \int_{0}^{1}G\Psi_n d\eta
\end{gather*}

The transformed equation is obtained and admits an analytical solution.
\begin{columns}[t]
\begin{column}{5cm}
	For $\lambda > 0$:
	\small
\begin{gather*}
%\text{For}\quad\lambda > 0\rightarrow
\bar{\Theta}_n''+(-\beta^2\lambda_n^2-\Bi\gamma)\bar{\Theta}_n=-\bar{G}_n(\xi)\\
\bar{\Theta}_n'(0)=0; \quad \bar{\Theta}_n'(1)=0\\
%\end{gather*}
%\begin{gather*}
\text{where} \\\bar{G}_n(\xi)= \int_{0}^{1}G(\xi,\eta)\Psi_n(\eta)\dd\eta
\end{gather*}
\end{column}
\begin{column}{5cm}
	 For $\lambda = 0$:
	 \small
	 \begin{gather*}
	 \bar{\Theta}_0''-(\Bi\gamma)\bar{\Theta}_0=-\bar{G}_0(\xi)\\
	 \bar{\Theta}_0'(0)=0; \quad \bar{\Theta}_0'(1)=0\\
	 %\end{gather*}
	 %\begin{gather*}
	 \text{where}\\\bar{G}_0(\xi)= \int_{0}^{1}G(\xi,\eta)\dd\eta
	 \end{gather*}
\end{column}
\end{columns}
\end{frame}
%^^^^^^^^^^^^^^^^^^^^^^^^^^
%^^^^^^^^^^^^^^^^^^^^^^^^^^




		 
		 

		 
		 
		 
		 
		 
%%%%%%%%%%%%%%%%%%%%%%%%%%%%
%%%%%%%%%%%%%%%%%%%%%%%%%%%%	
\begin{frame}{Classical Integral Transform Technique}\includelogo
			
The transformed equations admit a analytical solutions. \\[1em]
	
In order to obtain the final temperature field, the inversion formula must be utilized and the summation must be truncated to a finite value ($\nmax$).
		
\begin{align*}
\text{Inversion} \quad & \Rightarrow \quad \Theta= \sum_{n=0}^{\nmax} \frac{ \bar{\Theta}_n(\xi)\Psi_n(\eta)}{N_n} 
\end{align*}		
\end{frame}
%^^^^^^^^^^^^^^^^^^^^^^^^^^
%^^^^^^^^^^^^^^^^^^^^^^^^^^




		








%~~~~~~~~~~~~~~~~~~~~~~~~~~		
\section{Results}
%~~~~~~~~~~~~~~~~~~~~~~~~~~






%%%%%%%%%%%%%%%%%%%%%%%%%%%%
%%%%%%%%%%%%%%%%%%%%%%%%%%%%
\begin{frame}{Heat Generation}\includelogo
\scriptsize{
	\begin{gather*}\label{eq:heatgen}
	%G(\xi,\eta)=\frac{Gen}{\sigma^2\pi}\exp\left[-\frac{{(\xi-0.25)^2+(\eta-0.4)^2}}{\sigma^2}\right]
	G(\xi,\eta)=\frac{G_e}{\sigma^2\pi}\exp\left[-\frac{{(\xi-\xi_0)^2+(\eta-\eta_0)^2}}{\sigma^2}\right]
	\end{gather*}
For this current work, two cases are considered:
\begin{itemize}
	\item Centered HG where $\xi_0=1/2$ and $\eta_0=1/2$
	\item Off-Centered HG where $\xi_0=1/4$ and $\eta_0=2/5$
\end{itemize} 
For both cases $G_e=1$ and $\sigma=0.1$. }
\vspace*{-1.5em}
\begin{figure}
	\centering
	%\caption{Contour plot of the heat generation considered cases.} \label{heatgenimage}
	%\hspace*{-6cm}
	\subfloat[][Centered HG]{%: $\xi_0=1/2$ and $\eta_0=1/2$]{%\label{centralheat}
		\fbox{\includegraphics[width=0.25\hsize]{Figuras/contour-geracao-central.pdf}}}
	\hspace*{1cm}
	\subfloat[][ Off-Centered HG]{%  $\xi_0=1/4$ and $\eta_0=2/5$]{%\label{initialheat}
		\fbox{\includegraphics[width=0.25\hsize]{Figuras/contour-geracao-assimetrica.pdf}}}
	%\\
	%\source{Citação da fonte ou `O autor'. (opcional)}
\end{figure}
\end{frame}
%^^^^^^^^^^^^^^^^^^^^^^^^^^
%^^^^^^^^^^^^^^^^^^^^^^^^^^











%%%%%%%%%%%%%%%%%%%%%%%%%%%%
%%%%%%%%%%%%%%%%%%%%%%%%%%%%
\begin{frame}{Results}\includelogo
\tiny
\begin{table}[h!]
\caption{$\Theta(\xi,\eta)$ convergence solved for Centered HG by CITT}\vspace{-.7em} %\label{tab:geracaonomeio}
\centering
\setlength\tabcolsep{1.5pt}
\renewcommand\arraystretch{1}
	\begin{tabular}{|c|c|c|c|c|c|c|c|c|}%c|
		\hline
		&\multicolumn{4}{|c|}{$\Bi\gamma$=1 \& $\beta$=1}& \multicolumn{4}{c|}{$\Bi\gamma$=0.1 \& $\beta$=1} \\ \hline
		$\nmax$ &$\Theta$(0.2,0.4)& $\Theta$(0.5,0.5)&  $\Theta$(0.6, 0.8)& $\Theta$(0.9,0.9)&$\Theta$(0.2, 0.4)& $\Theta$(0.5, 0.5)&$\Theta$(0.6,0.8)& $\Theta$(0.9,0.9) \\\hline
		1 &          0.981222 & 1.05643 & 1.03739 & 0.966733  & 9.98087 & 10.0575 & 10.0382 & 9.96593 \\
		5 &          1.00221 & 1.18948 & 0.997708 & 0.954075  &  \bfseries 10.0025 & 10.1922 & 9.99803 &   \bfseries 9.95274 \\
		10 & \bfseries 1.00213 & 1.20298 & 1.00209 &  \bfseries 0.954083  & 10.0025 & 10.2058 & 10.0024 & 9.95274 \\
		%15 & 1.00213 & 1.20322 & 1.00213 & 0.954083 & 15 & 10.0025 & 10.2060 & 10.0025 & 9.95274 \\
		20 & 1.00213 &  \bfseries 1.20323 &  \bfseries 1.00213 & 0.954083  & 10.0025 &  \bfseries 10.2060 &  \bfseries  10.0025 & 9.95274 \\
		%25 & 1.00213 & 1.20323 & 1.00213 & 0.954083 & 25 & 10.0025 & 10.2060 & 10.0025 & 9.95274 \\
		30 & 1.00213 & 1.20323 & 1.00213 & 0.954083  & 10.0026 & 10.2060 & 10.0025 & 9.95274 \\
		\hline
		&\multicolumn{4}{|c|}{$\Bi\gamma$=1 \& $\beta$=0.8} & \multicolumn{4}{c|}{$\Bi\gamma$=0.1 \& $\beta$=0.8} \\ \hline
		$\nmax$ & $\Theta$(0.2,0.4)& $\Theta$(0.5,0.5)&  $\Theta$(0.6,0.8)& $\Theta$(0.9,0.9) & $\Theta$(0.2,0.4)& $\Theta$(0.5,0.5)&$\Theta$(0.6,0.8)& $\Theta$(0.9,0.9)\\ \hline
		1 & 0.981222 & 1.05643 & 1.03739 & 0.966733 & 9.98087 & 10.0575 & 10.0382 & 9.96593 \\
		5 & 1.01980 & 1.23495 & 0.980675 & 0.940388 & 10.0212 & 10.2395 & 9.98041 & 9.93805 \\
		10 & \bfseries  1.01957 & 1.25429 & 0.987398 & \bfseries  0.940418  &  \bfseries 10.0210 & 10.2589 & 9.98715 &  \bfseries 9.93808 \\
		%15 & 1.01957 & 1.25466 & 0.987466 & 0.940418 & 15 & 10.0210 & 10.2593 & 9.98722 & 9.93808 \\
		20 & 1.01957 &  \bfseries 1.25467 &  \bfseries 0.987461 & 0.940418 & 10.0210 &  \bfseries 10.2593 & \bfseries  9.98722 & 9.93808 \\
		%25 & 1.01957 & 1.25467 & 0.987461 & 0.940418 & 25 & 10.0210 & 10.2593 & 9.98722 & 9.93808 \\
		30 & 1.01957 & 1.25467 & 0.987461 & 0.940418 & 10.0210 & 10.2593 & 9.98722 & 9.93808 \\
		\hline
		&\multicolumn{4}{|c|}{$\Bi\gamma$=1 \& $\beta$=0.5} & \multicolumn{4}{c|}{$\Bi\gamma$=0.1 \& $\beta$=0.5}\\ \hline
		$\nmax$ & $\Theta$(0.2,0.4)& $\Theta$(0.5,0.5)&  $\Theta$(0.6,0.8)& $\Theta$(0.9,0.9)  & $\Theta$(0.2,0.4)& $\Theta$(0.5,0.5)&  $\Theta$(0.6,0.8)& $\Theta$(0.9,0.9)\\\hline
		1 & 0.981222 & 1.05643 & 1.03739 & 0.966733 & 9.98087 & 10.0575 & 10.0382 & 9.96593 \\
		5 & 1.09885 & 1.38536 & 0.921228 & 0.874113  & 10.1105 & 10.4039 & 9.91626 & 9.86204 \\
		10 & 1.09731 & 1.42441 & 0.936699 & 0.874431  &  \bfseries 10.1090 & 10.4432 & 9.93187 &  \bfseries 9.86236 \\
		%15 & 1.09730 & 1.42523 & 0.936875 & 0.874429 & 15 & 10.1090 & 10.4440 & 9.93204 & 9.86236 \\
		20 &  \bfseries 1.09730 &  \bfseries 1.42526 &  \bfseries 0.936864 &  \bfseries 0.874429 & 10.1090 &  \bfseries 10.4441 &  \bfseries 9.93203 & 9.86236 \\
		%25 & 1.09730 & 1.42526 & 0.936864 & 0.874429 & 25 & 10.1090 & 10.4441 & 9.93203 & 9.86236 \\
		30 & 1.09730 & 1.42526 & 0.936864 & 0.874429 & 10.1090 & 10.4441 & 9.93203 & 9.86236 \\
		\hline
	\end{tabular}
\end{table} 
\end{frame}
%^^^^^^^^^^^^^^^^^^^^^^^^^^
%^^^^^^^^^^^^^^^^^^^^^^^^^^












%%%%%%%%%%%%%%%%%%%%%%%%%%%%
%%%%%%%%%%%%%%%%%%%%%%%%%%%%
\begin{frame}{Results}\includelogo
\tiny
\begin{table}[h!]
	\caption{$\Theta(\xi,\eta)$ convergence solved for Off-Centered HG by CITT}\vspace{-.7em} %\label{tab:geracaodeslocada}
	\centering
	\setlength\tabcolsep{1.5pt}
	\renewcommand\arraystretch{1}
	\begin{tabular}{|c|c|c|c|c|c|c|c|c|}%c|
		\hline
		&\multicolumn{4}{|c|}{$\Bi\gamma$=1 \& $\beta$=1}& \multicolumn{4}{c|}{$\Bi\gamma$=0.1 \& $\beta$=1} \\ \hline
		$\nmax$ & $\Theta$(0.2,0.4)& $\Theta$(0.5,0.5)&  $\Theta$(0.6,0.8)& $\Theta$(0.9,0.9) & $\Theta$(0.2,0.4)& $\Theta$(0.5,0.5)&  $\Theta$(0.6,0.8)& $\Theta$(0.9,0.9) \\\hline
		1 & 1.14595 & 0.991937 & 0.920318 & 0.865087  & 10.1555 & 9.99006 & 9.91129 & 9.84928 \\
		5 & 1.26434 & 1.02009 & 0.914914 &  \bfseries 0.863013 & 10.2754 & 10.0190 & 9.90570 &  \bfseries 9.84707  \\
		10 & 1.27061 &  \bfseries 1.01987 &  \bfseries 0.91489 & 0.863013 & 10.2817 &  \bfseries 10.0188 &  \bfseries 9.90568 & 9.84707 \\
		%15 &  1.27082 & 1.01987 & 0.91489 & 0.863013 & 15 & 10.2819 & 10.0188 & 9.90568 & 9.84707 \\
		20 &  \bfseries 1.27083 & 1.01987 & 0.91489 & 0.863013 &  \bfseries 10.2819 & 10.0188 & 9.90568 & 9.84707 \\
		%25 &  1.27083 & 1.01987 & 0.91489 & 0.863013 & 25 & 10.2819 & 10.0188 & 9.90568 & 9.84707 \\
		30 & 1.27083 & 1.01987 & 0.91489 & 0.863013 & 10.2819 & 10.0188 & 9.90568 & 9.84707 \\
		\hline
		&\multicolumn{4}{|c|}{$\Bi\gamma$=1 \& $\beta$=0.8} & \multicolumn{4}{c|}{$\Bi\gamma$=0.1 \& $\beta$=0.8} \\ \hline
		$\nmax$ & $\Theta$(0.2,0.4)& $\Theta$(0.5,0.5)&  $\Theta$(0.6,0.8)& $\Theta$(0.9,0.9) & $\Theta$(0.2,0.4)& $\Theta$(0.5,0.5)&  $\Theta$(0.6,0.8)& $\Theta$(0.9,0.9)\\ \hline
		1 &  1.1564 & 0.991937 & 0.900772 & 0.846859 & 10.1682 & 9.99006 & 9.88666 & 9.82559 \\
		5 & 1.32297 & 1.04002 & 0.890144 & \bfseries  0.840922 & 10.3380 & 10.0401 & 9.87552 &  \bfseries 9.81914 \\
		10 & 1.33221 &  \bfseries 1.03950 &  \bfseries 0.890058 & 0.840922 & 10.3473 &  \bfseries 10.0396 &  \bfseries 9.87544 & 9.81914 \\
		%15 & 1.33253 & 1.03950 & 0.890058 & 0.840922 & 15 & 10.3476 & 10.0396 & 9.87544 & 9.81914 \\
		20 &  \bfseries 1.33253 & 1.03950 & 0.890058 & 0.840922 &  \bfseries 10.3476 & 10.0396 & 9.87544 & 9.81914 \\
		%25 & 1.33253 & 1.03950 & 0.890058 & 0.840922 & 25 & 10.3476 & 10.0396 & 9.87544 & 9.81914 \\
		30 & 1.33253 & 1.03950 & 0.890058 & 0.840922 & 10.3476 & 10.0396 & 9.87544 & 9.81914 \\
		\hline
		&\multicolumn{4}{|c|}{$\Bi\gamma$=1 \& $\beta$=0.5} & \multicolumn{4}{c|}{$\Bi\gamma$=0.1 \& $\beta$=0.5}\\ \hline
		$\nmax$ & $\Theta$(0.2,0.4)& $\Theta$(0.5,0.5)&  $\Theta$(0.6,0.8)& $\Theta$(0.9,0.9) & $\Theta$(0.2,0.4)& $\Theta$(0.5,0.5)&  $\Theta$(0.6,0.8)& $\Theta$(0.9, 0.9)\\\hline
		1 & 1.18823 & 0.991937 & 0.830465 & 0.772371 & 10.2157 & 9.99006 & 9.77686 & 9.70579 \\
		5 &  1.52152 & 1.12367 & 0.794684 & 0.738850 & 10.5654 & 10.1341 & 9.73726 & 9.66629 \\
		10 &  1.54154 & 1.12104 & 0.793933 &  \bfseries 0.738828 & 10.5855 & 10.1315 & 9.73650 &  \bfseries 9.66627 \\
		%15 & 1.54227 & 1.12102 & 0.793928 & 0.738828 & 15 & 10.5862 & 10.1314 & 9.73649 & 9.66627 \\
		20 & 1.54228 &  \bfseries 1.12102 &  \bfseries 0.793928 & 0.738828 &  \bfseries 10.5862 &  \bfseries 10.1314 &  \bfseries 9.73649 & 9.66627 \\
		%25 & 1.54229 & 1.12102 & 0.793928 & 0.738828 & 25 & 10.5862 & 10.1314 & 9.73649 & 9.66627  \\
		30 &  \bfseries 1.54229 & 1.12102 & 0.793928 & 0.738828 & 10.5862 & 10.1314 & 9.73649 & 9.66627 \\
		\hline
	\end{tabular}
\end{table} 
\end{frame}
%^^^^^^^^^^^^^^^^^^^^^^^^^^
%^^^^^^^^^^^^^^^^^^^^^^^^^^












%%%%%%%%%%%%%%%%%%%%%%%%%%%%
%%%%%%%%%%%%%%%%%%%%%%%%%%%%
\begin{frame}{Results}\includelogo
\begin{figure}[!h]
	\centering
	\caption{Illustration of the solution by CITT with the different heat generation positions with the parameters of $\Bi\gamma$=1 and $\beta$=1.} 
	\subfloat[][Off-Centered HG case ]{
		\fbox{\includegraphics[width=0.30\hsize]{Figuras/contour-citt-assimetrico.pdf}}}
	\hspace*{1cm} %\\
	\subfloat[][Centered HG case]{
		\fbox{\includegraphics[width=0.30\hsize]{Figuras/contour-citt-central.pdf}}}
\end{figure}
\end{frame}
%^^^^^^^^^^^^^^^^^^^^^^^^^^
%^^^^^^^^^^^^^^^^^^^^^^^^^^	






%~~~~~~~~~~~~~~~~~~~~~~~~~~
\section{Conclusions}
%~~~~~~~~~~~~~~~~~~~~~~~~~~






%%%%%%%%%%%%%%%%%%%%%%%%%%%%
%%%%%%%%%%%%%%%%%%%%%%%%%%%%
\begin{frame}{Conclusions}\includelogo
\begin{itemize}
	 \item The solution is fully analytical
	 \vspace{2em}
	 \item CITT had a great performance to obtain high accuracy with very few terms in the solution summation.
	 \vspace{2em}
	 \item This is a preliminary work and for near future works, one can propose: Use of the complete Heat Generation and variation of the heat conductivity. 
 \end{itemize}
\end{frame}
%^^^^^^^^^^^^^^^^^^^^^^^^^^
%^^^^^^^^^^^^^^^^^^^^^^^^^^





	
	

	
	
%%%%%%%%%%%%%%%%%%%%%%%%%%%%
%%%%%%%%%%%%%%%%%%%%%%%%%%%%	
\begin{frame}{Acknowledgments}\includelogo
	\begin{center}
	\Huge   \textbf{Thank You}
	\end{center}
	
	\vspace{3em}
	\hrule
	\vspace{2mm}
	
	\textbf{Acknowledgments:}
	\vspace{2mm}
	
    \begin{center}
	\includegraphics[height=2em]{Agencies/Capes.jpg}
	\hspace{5mm}
	\includegraphics[height=2em]{Agencies/Cnpq.png}
	\hspace{5mm}
	\includegraphics[height=2em]{Agencies/Faperj.jpg}
	\hspace{5mm}\\[5mm] 
		\includegraphics[height=4em]{LogoPPGEM-color.png}
	\hspace{5mm}
	\includegraphics[height=4em]{Logo-UERJ.png}
	\hspace{0mm}

    \end{center}
\end{frame}
%^^^^^^^^^^^^^^^^^^^^^^^^^^
%^^^^^^^^^^^^^^^^^^^^^^^^^^
	







	
%%%%%%%%%%%%%%%%%%%%%%%%%%%%
%%%%%%%%%%%%%%%%%%%%%%%%%%%%
\begin{frame}\includelogo
\begin{columns}
\column{0.5\textwidth}
\frametitle{Lists - Examples}
\begin{itemize}
\item Point A
\item Point B
\begin{itemize}
\item part 1
\item part 2
\end{itemize}
\item Point C
\end{itemize}

\begin{enumerate}[I]
\item Point A
\item Point B
\begin{enumerate}[i]
\item part 1
\item part 2
\end{enumerate}
\item Point C
\end{enumerate}
\column{0.5\textwidth}
\begin{description}
\item[API] Application Programming Interface
\item[LAN] Local Area Network
\item[ASCII] American Standard Code for Information Interchange
\end{description}

\end{columns}
\end{frame}
%^^^^^^^^^^^^^^^^^^^^^^^^^^
%^^^^^^^^^^^^^^^^^^^^^^^^^^









%%%%%%%%%%%%%%%%%%%%%%%%%%%%
%%%%%%%%%%%%%%%%%%%%%%%%%%%%
\begin{frame}{Block Texts  - Examples}\includelogo

\begin{block}{Block Title}
Lorem ipsum dolor sit amet, consectetur adipisicing elit, 
sed do eiusmod 
\end{block}

\begin{alertblock}{Block Title}
Lorem ipsum dolor sit amet, consectetur adipisicing elit, 
sed do eiusmod tempor
\end{alertblock}

\begin{definition}
A prime number is a number that...
\end{definition}

\begin{example}
Lorem ipsum dolor sit amet, consectetur adipisicing elit, 
sed do eiusmod tempor incididunt ut labore et
dolore magna aliqua.
\end{example}
    
\end{frame}
%^^^^^^^^^^^^^^^^^^^^^^^^^^
%^^^^^^^^^^^^^^^^^^^^^^^^^^








%%%%%%%%%%%%%%%%%%%%%%%%%%%%
%%%%%%%%%%%%%%%%%%%%%%%%%%%%
\begin{frame}[fragile]{Including Code  - Examples}\includelogo
\begin{lstlisting}
\begin{gather}
\frac{\partial^2 T}{\partial x^2} + 
\frac{\partial^2 T}{\partial y^2} 
 = 0
\end{gather}
\end{lstlisting}
\end{frame}
%^^^^^^^^^^^^^^^^^^^^^^^^^^
%^^^^^^^^^^^^^^^^^^^^^^^^^^









%%%%%%%%%%%%%%%%%%%%%%%%%%%%
%%%%%%%%%%%%%%%%%%%%%%%%%%%%
%\bibliographystyle{plain}
%	\begin{frame}{References}\includelogo
%		\bibliography{livia}
%	\end{frame}
%^^^^^^^^^^^^^^^^^^^^^^^^^^
%^^^^^^^^^^^^^^^^^^^^^^^^^^



\end{document}